\exnumber{06}
\extitle{UX IV - Ex 7 - Test}
\exfiles{ex06/archive zip, nomme Name\_FirstName\_DeliverableName\_Date\_VersionNumber}
\exauthorize{Those named in the exercise}
\makeheaderfiles

\textbf{Instructions :}

Test your wireframe prototype with 5 different people and write down the 5 main iterations that need to be made on the next version.

Cet exercice se fait seul.

\textbf{Livrables :}

\begin{itemize}
\item A PDF document with the usability test script
\item Des photos des tests en cours
\item A PDF docuemnt with a report on what you learnt and a list of the 5 main iterations that you'd like to do on the next version
\item Déposez les documents sur Github, dans un zip nommé « ProjectTitle\_Name\_FirstName », et chaque livrable ainsi : « Name\_FirstName\_DeliverableName\_Date\_VersionNumber ». Exemple : « Doe\_John\_MidFiPrototype\_05242024\_V1 ».
\item Ajoutez un titre dans tous les documents écrits.
\item Soignez l'aspect global des livrables. Ils doivent être clairs, simples et lisibles. Vous pouvez regarder canva.com, sans surcharger le dessin de détails.
\end{itemize}

\textbf{Ressources :}

\begin{itemize}
\item Running a usability test \url{https://www.usability.gov/how-to-and-tools/methods/running-usability-tests.html}
\item Usability testing: what is it and how to do it? \url{https://uxdesign.cc/usability-testing-what-is-it-how-to-do-it-51356e5de5d}
\item Analyzing usability testing data \url{https://uxdesign.cc/analysing-usability-testing-data-97667ae4999e}
\item The art of guerrilla usability testing \url{http://www.uxbooth.com/articles/the-art-of-guerrilla-usability-testing/}
\item How to write a user testing report that people will actually read \url{https://uxdesign.cc/how-to-write-a-user-testing-report-that-people-will-actually-read-652d15d2f92e}
\end{itemize}

\textbf{Citation :}

Prototypes are created for testing purposes, that's why creating a prototype without running tests on it makes no sense.

